\chapter{Input/Output}
\label{chap:io}

Input and output on the Amiga spans two distinct worlds: the CLI (Command Line Interface) where text-based programs read arguments and write output, and Workbench where programs are launched from icons with tooltypes. Novus provides unified support for both, with RAII wrappers that prevent the resource leaks that plagued classic Amiga development.

%%%%%%%%%%%%%%%%%%%%%%%%%%%%%%%%%%%%%%%%%%%%%%%%%%%%%%%%%%%%%%%%%%%%%%%%%%%%%%%
\section{Console Output}
\label{sec:console-output}
%%%%%%%%%%%%%%%%%%%%%%%%%%%%%%%%%%%%%%%%%%%%%%%%%%%%%%%%%%%%%%%%%%%%%%%%%%%%%%%

The \texttt{std::io::file} module provides basic console I/O functions:

\begin{lstlisting}
from std::io::file import write, print, println

fn hello() {
    // Simple string output
    print("Hello, ")
    println("Amiga!")  // Adds newline

    // Printf-style formatting (AmigaOS RawDoFmt syntax)
    let count: i32 = 42
    let name: *u8 = "Player"
    write("Score: %ld for %s\n", count, name)
}
\end{lstlisting}

\begin{warningbox}
\textbf{Use \texttt{\%ld} for 32-bit integers, not \texttt{\%d}!} AmigaOS \texttt{RawDoFmt} expects 16-bit values for \texttt{\%d}. Since Novus passes all integers as 32-bit, always use the \texttt{l} (long) prefix: \texttt{\%ld}, \texttt{\%lu}, \texttt{\%lx}.
\end{warningbox}

\subsection{Format Specifiers}

\begin{table}[h]
\centering
\begin{tabular}{ll}
\textbf{Specifier} & \textbf{Meaning} \\
\hline
\texttt{\%ld} & Signed 32-bit decimal \\
\texttt{\%lu} & Unsigned 32-bit decimal \\
\texttt{\%lx} & Hexadecimal (lowercase) \\
\texttt{\%lX} & Hexadecimal (uppercase) \\
\texttt{\%s} & Null-terminated string \\
\texttt{\%c} & Single character \\
\texttt{\%b} & BSTR (BPTR to byte-counted string) \\
\end{tabular}
\caption{AmigaOS format specifiers}
\end{table}

\subsection{Array-Based Formatting}

When building arguments programmatically:

\begin{lstlisting}
from std::io::file import write_array

fn formatted_output(major: i32, minor: i32, patch: i32) {
    var args: [i32; 3] = [major, minor, patch]
    write_array("Version: %ld.%ld.%ld\n", &args[0])
}
\end{lstlisting}

%%%%%%%%%%%%%%%%%%%%%%%%%%%%%%%%%%%%%%%%%%%%%%%%%%%%%%%%%%%%%%%%%%%%%%%%%%%%%%%
\section{Console Input}
\label{sec:console-input}
%%%%%%%%%%%%%%%%%%%%%%%%%%%%%%%%%%%%%%%%%%%%%%%%%%%%%%%%%%%%%%%%%%%%%%%%%%%%%%%

\begin{lstlisting}
from std::io::file import wait_for_key, key_pressed, read_key

fn input_examples() {
    // Block until any key pressed
    print("Press any key to continue...")
    wait_for_key()
    println("")

    // Non-blocking check
    if key_pressed() {
        let ch = read_key()
        write("You pressed: %c\n", (i32)ch)
    }

    // Blocking read of single key
    print("Enter a character: ")
    let input = read_key()  // Returns u8
    write("Got: %c\n", (i32)input)  // Cast to i32 for format string
}
\end{lstlisting}

\begin{notebox}
\texttt{read\_key()} returns \texttt{u8}. The cast to \texttt{i32} is required because AmigaOS \texttt{\%c} format expects a 32-bit value on the stack.
\end{notebox}

%%%%%%%%%%%%%%%%%%%%%%%%%%%%%%%%%%%%%%%%%%%%%%%%%%%%%%%%%%%%%%%%%%%%%%%%%%%%%%%
\section{ANSI Console Control}
\label{sec:ansi}
%%%%%%%%%%%%%%%%%%%%%%%%%%%%%%%%%%%%%%%%%%%%%%%%%%%%%%%%%%%%%%%%%%%%%%%%%%%%%%%

The Amiga console supports ANSI X3.64 escape sequences for cursor control, colors, and text styling. The \texttt{std::io::ansi} module provides type-safe access to these sequences.

\begin{notebox}
\textbf{Return types:} Simple functions like \texttt{reset()}, \texttt{bold()}, \texttt{clear\_screen()} return \texttt{Str} (static escape sequences). Parameterized functions like \texttt{move\_to()}, \texttt{fg()}, \texttt{cursor\_up()} return \texttt{String} (dynamically built). Both have \texttt{as\_cstr()} for use with \texttt{write()}.
\end{notebox}

\subsection{Basic Functions}

\begin{lstlisting}
from std::io::ansi import clear_screen, home, reset, bold, underline, reverse
from std::io::file import write

fn styled_output() {
    // Clear screen and move to top-left
    write("%s%s", clear_screen().as_cstr(), home().as_cstr())

    // Styled text
    write("%sThis is bold%s\n", bold().as_cstr(), reset().as_cstr())
    write("%sThis is underlined%s\n", underline().as_cstr(), reset().as_cstr())
    write("%sThis is reversed%s\n", reverse().as_cstr(), reset().as_cstr())
}
\end{lstlisting}

\subsection{Pen Colors}

The Amiga console uses \emph{pen numbers}, not fixed colors. The actual color depends on the Workbench palette:

\begin{lstlisting}
from std::io::ansi import Pen, fg, bg, reset
from std::io::file import write

fn colored_output() {
    // Foreground colors
    write("%sText in Pen 1%s\n", fg(Pen::Pen1).as_cstr(), reset().as_cstr())
    write("%sText in Pen 2%s\n", fg(Pen::Pen2).as_cstr(), reset().as_cstr())
    write("%sText in Pen 3%s\n", fg(Pen::Pen3).as_cstr(), reset().as_cstr())

    // Background colors
    write("%sHighlighted%s\n", bg(Pen::Pen3).as_cstr(), reset().as_cstr())
}
\end{lstlisting}

\begin{table}[h]
\centering
\begin{tabular}{lll}
\textbf{Pen} & \textbf{Workbench 2.0+ Default} & \textbf{ANSI Code} \\
\hline
Pen0 & Grey (background) & 30/40 \\
Pen1 & Black (text) & 31/41 \\
Pen2 & White (highlight) & 32/42 \\
Pen3 & Blue (shine/fill) & 33/43 \\
Pen4--7 & Only on 8+ color screens & 34--37/44--47 \\
\end{tabular}
\caption{Amiga console pen mapping}
\end{table}

\begin{notebox}
On 4-color screens, pens 4--7 are remapped: Pen4 becomes invisible (same as background), Pen5 maps to Pen1, Pen6 to Pen2, Pen7 to Pen3.
\end{notebox}

\subsection{Cursor Movement}

\begin{lstlisting}
from std::io::ansi import move_to, cursor_up, cursor_down, cursor_left, cursor_right
from std::io::file import write

fn cursor_demo() {
    // Absolute positioning (1-indexed, row then column)
    write("%s", move_to(10, 20).as_cstr())
    write("At row 10, column 20")

    // Relative movement
    write("%s", cursor_up(5).as_cstr())    // Move up 5 rows
    write("%s", cursor_right(10).as_cstr()) // Move right 10 columns
}
\end{lstlisting}

\subsection{Screen and Line Clearing}

\begin{lstlisting}
from std::io::ansi import clear_screen, clear_line, clear_to_eol, clear_to_eos
from std::io::ansi import ClearMode, clear
from std::io::file import write

fn clearing_demo() {
    // Simple functions
    write("%s", clear_screen().as_cstr())  // Clear entire screen
    write("%s", clear_line().as_cstr())    // Clear current line
    write("%s", clear_to_eol().as_cstr())  // Clear to end of line
    write("%s", clear_to_eos().as_cstr())  // Clear to end of screen

    // Parameterized clearing
    write("%s", clear(ClearMode::ToEnd).as_cstr())   // From cursor to end
    write("%s", clear(ClearMode::ToStart).as_cstr()) // From start to cursor
    write("%s", clear(ClearMode::All).as_cstr())     // Entire screen
}
\end{lstlisting}

\subsection{The Ansi Builder}

For complex styling, use the builder pattern:

\begin{lstlisting}
from std::io::ansi import Ansi, Pen, reset
from std::io::file import write

fn builder_demo() {
    // Build complex style
    var style = Ansi::new()
    style.bold()
    style.fg(Pen::Pen2)
    style.bg(Pen::Pen1)

    // Apply and reset
    let seq = style.build()
    write("%sStyled Text%s\n", seq.as_cstr(), reset().as_cstr())
}

// Method chaining
fn chained_demo() {
    var style = Ansi::new()
    style.bold().underline().fg(Pen::Pen3)

    write("%sChained%s\n", style.build().as_cstr(), reset().as_cstr())
}
\end{lstlisting}

%%%%%%%%%%%%%%%%%%%%%%%%%%%%%%%%%%%%%%%%%%%%%%%%%%%%%%%%%%%%%%%%%%%%%%%%%%%%%%%
\section{File I/O}
\label{sec:file-io}
%%%%%%%%%%%%%%%%%%%%%%%%%%%%%%%%%%%%%%%%%%%%%%%%%%%%%%%%%%%%%%%%%%%%%%%%%%%%%%%

The \texttt{amiga::sys::dos} module provides RAII wrappers for AmigaOS file operations.

\subsection{OwnedFileHandle --- RAII File Access}

\begin{lstlisting}
from amiga::sys::dos import OwnedFileHandle
from amiga::raw::consts import MODE_OLDFILE, MODE_NEWFILE

fn read_file_example() -> Result<(), DosError> {
    // Open file - automatically closed when handle goes out of scope
    let file = OwnedFileHandle::open("S:Startup-Sequence", MODE_OLDFILE)?

    // Read content
    var buffer: [u8; 1024] = [0; 1024]
    let bytes_read = file.read((*u8)&buffer[0], 1024)

    if bytes_read > 0 {
        // Process buffer...
    }

    return Result::Ok(())
}  // file automatically closed here!
\end{lstlisting}

\begin{comingfromc}
In C, you'd use \texttt{Open()} and must remember to call \texttt{Close()} on every code path---including error paths. With \texttt{OwnedFileHandle}, the file is closed automatically when the handle goes out of scope, even if an error occurs.
\end{comingfromc}

\subsection{File Open Modes}

\begin{table}[h]
\centering
\begin{tabular}{lll}
\textbf{Mode} & \textbf{Value} & \textbf{Description} \\
\hline
\texttt{MODE\_OLDFILE} & 1005 & Open existing file (read/write) \\
\texttt{MODE\_NEWFILE} & 1006 & Create new file (truncate if exists) \\
\texttt{MODE\_READWRITE} & 1004 & Open with shared lock \\
\end{tabular}
\caption{DOS file open modes}
\end{table}

\subsection{Reading and Writing}

\begin{lstlisting}
from amiga::sys::dos import OwnedFileHandle
from amiga::raw::consts import MODE_NEWFILE

fn write_file_example() -> Result<(), DosError> {
    let file = OwnedFileHandle::open("RAM:output.txt", MODE_NEWFILE)?

    // Write bytes
    let data: *u8 = "Hello, Amiga!\n"
    let written = file.write(data, 14)

    // Write all bytes (loops until complete)
    file.write_all(data, 14)?

    return Result::Ok(())
}

fn read_all_example() -> Result<(), DosError> {
    let file = OwnedFileHandle::open("RAM:input.txt", MODE_OLDFILE)?

    var buffer: [u8; 4096] = [0; 4096]

    // Read up to buffer size
    let total = file.read_all((*u8)&buffer[0], 4096)?

    write("Read %ld bytes\n", (i32)total)

    return Result::Ok(())
}
\end{lstlisting}

\subsection{Reading Lines}

\begin{lstlisting}
from amiga::sys::dos import OwnedFileHandle
from amiga::raw::consts import MODE_OLDFILE

fn read_lines_example() -> Result<(), DosError> {
    let file = OwnedFileHandle::open("S:Startup-Sequence", MODE_OLDFILE)?

    var line_buffer: [u8; 256] = [0; 256]
    var line_num: i32 = 1

    loop {
        match file.read_line((*u8)&line_buffer[0], 256) {
            Result::Ok(len) => {
                if len == 0 {
                    break  // EOF
                }
                write("%ld: %s", line_num, (*u8)&line_buffer[0])
                line_num++
            },
            Result::Err(_) => break,
        }
    }

    return Result::Ok(())
}
\end{lstlisting}

\subsection{Seeking}

\begin{lstlisting}
from amiga::sys::dos import OwnedFileHandle
from amiga::raw::consts import MODE_OLDFILE, OFFSET_BEGINNING, OFFSET_CURRENT, OFFSET_END

fn seek_example() -> Result<(), DosError> {
    let file = OwnedFileHandle::open("RAM:data.bin", MODE_OLDFILE)?

    // Seek to beginning
    file.seek(0, OFFSET_BEGINNING)

    // Seek forward 100 bytes from current position
    file.seek(100, OFFSET_CURRENT)

    // Get file size by seeking to end
    let size = file.seek(0, OFFSET_END)
    write("File size: %ld bytes\n", size)

    // Seek back to start
    file.seek(0, OFFSET_BEGINNING)

    return Result::Ok(())
}
\end{lstlisting}

\begin{table}[h]
\centering
\begin{tabular}{lll}
\textbf{Constant} & \textbf{Value} & \textbf{Meaning} \\
\hline
\texttt{OFFSET\_BEGINNING} & $-1$ & From start of file \\
\texttt{OFFSET\_CURRENT} & 0 & From current position \\
\texttt{OFFSET\_END} & 1 & From end of file \\
\end{tabular}
\caption{DOS seek offset constants}
\end{table}

\subsection{Ownership Transfer}

When you need to pass a file handle to external code:

\begin{lstlisting}
fn ownership_example() -> Result<i32, DosError> {
    let file = OwnedFileHandle::open("RAM:test.txt", MODE_OLDFILE)?

    // Transfer ownership - prevents automatic close
    let raw_handle: i32 = file.into_raw()

    // raw_handle is now YOUR responsibility!
    // You must call Close(raw_handle) when done

    return Result::Ok(raw_handle)
}

fn from_raw_example(handle: i32) {
    // Take ownership of an existing handle
    let file = OwnedFileHandle::from_raw(handle)

    // Use file...

}  // Handle closed automatically
\end{lstlisting}

\begin{warningbox}
\textbf{Never use \texttt{from\_raw()} on system handles!} Handles like \texttt{Input()}, \texttt{Output()}, and \texttt{Error()} are owned by the system and must not be closed. Only use \texttt{from\_raw()} on handles you own.
\end{warningbox}

%%%%%%%%%%%%%%%%%%%%%%%%%%%%%%%%%%%%%%%%%%%%%%%%%%%%%%%%%%%%%%%%%%%%%%%%%%%%%%%
\section{Directory Locks}
\label{sec:dir-locks}
%%%%%%%%%%%%%%%%%%%%%%%%%%%%%%%%%%%%%%%%%%%%%%%%%%%%%%%%%%%%%%%%%%%%%%%%%%%%%%%

AmigaOS uses ``locks'' to access directories. \texttt{DirLockHandle} provides RAII management:

\begin{lstlisting}
from amiga::sys::dos import DirLockHandle
from amiga::raw::consts import ACCESS_READ, ACCESS_WRITE

fn lock_example() -> Result<(), DosError> {
    // Lock directory for reading
    let lock = DirLockHandle::lock("SYS:", ACCESS_READ)?

    // Get raw lock for FFI
    let raw_lock = lock.handle()

    // Use with Examine(), ExNext(), etc...

    return Result::Ok(())
}  // Lock automatically released
\end{lstlisting}

\begin{warningbox}
\textbf{Lock leaks prevent floppy ejection!} If a lock on a floppy disk isn't released, the user cannot eject the disk and sees ``Volume is in use.'' Always use \texttt{DirLockHandle} to ensure locks are released.
\end{warningbox}

%%%%%%%%%%%%%%%%%%%%%%%%%%%%%%%%%%%%%%%%%%%%%%%%%%%%%%%%%%%%%%%%%%%%%%%%%%%%%%%
\section{File System Helpers}
\label{sec:fs-helpers}
%%%%%%%%%%%%%%%%%%%%%%%%%%%%%%%%%%%%%%%%%%%%%%%%%%%%%%%%%%%%%%%%%%%%%%%%%%%%%%%

\begin{lstlisting}
from amiga::sys::dos import file_exists, delete_file, create_dir, file_info

fn fs_examples() -> Result<(), DosError> {
    // Check existence
    if file_exists("RAM:test.txt") {
        println("File exists!")
    }

    // Delete file
    delete_file("RAM:temp.txt")?

    // Create directory
    create_dir("RAM:NewFolder")?

    // Get file information
    match file_info("SYS:C/Dir") {
        Result::Ok(info) => {
            if info.is_file {
                write("Size: %ld bytes\n", info.size)
            } else if info.is_directory {
                println("It's a directory")
            }
        },
        Result::Err(_) => println("Cannot get info"),
    }

    return Result::Ok(())
}
\end{lstlisting}

The \texttt{FileInfo} structure:

\begin{lstlisting}
pub struct FileInfo {
    is_file: bool,      // True if regular file
    is_directory: bool, // True if directory
    size: i32,          // File size in bytes
    protection: i32,    // Protection bits (FIBF_*)
}
\end{lstlisting}

%%%%%%%%%%%%%%%%%%%%%%%%%%%%%%%%%%%%%%%%%%%%%%%%%%%%%%%%%%%%%%%%%%%%%%%%%%%%%%%
\section{Command-Line Arguments}
\label{sec:args}
%%%%%%%%%%%%%%%%%%%%%%%%%%%%%%%%%%%%%%%%%%%%%%%%%%%%%%%%%%%%%%%%%%%%%%%%%%%%%%%

AmigaOS uses \texttt{ReadArgs()} for parsing command-line arguments with a template syntax. The \texttt{amiga::workbench} module provides a type-safe wrapper.

\subsection{Basic Argument Parsing}

\begin{lstlisting}
from amiga::workbench import Args

pub fn main() -> i32 {
    match Args::parse("FROM/A,TO/A,VERBOSE/S") {
        Result::Ok(args) => {
            // Get required string arguments
            let from = args.get_str("FROM") or { return 10 }
            let to = args.get_str("TO") or { return 10 }

            // Get optional switch
            let verbose = args.get_switch("VERBOSE")

            if verbose {
                write("Copying %s to %s\n", from.as_cstr(), to.as_cstr())
            }

            // ... do work ...

            return 0
        },
        Result::Err(e) => {
            write("Error: %s\n", e.message())
            return 10
        }
    }
}
\end{lstlisting}

\subsection{Template Syntax}

\begin{table}[h]
\centering
\begin{tabular}{lll}
\textbf{Modifier} & \textbf{Meaning} & \textbf{Example} \\
\hline
(none) & Optional positional & \texttt{FILE} \\
\texttt{/A} & Required argument & \texttt{FILE/A} \\
\texttt{/S} & Switch (boolean flag) & \texttt{VERBOSE/S} \\
\texttt{/K} & Keyword (NAME=value) & \texttt{OUTPUT/K} \\
\texttt{/N} & Numeric argument & \texttt{COUNT/N} \\
\texttt{/M} & Multiple values & \texttt{FILES/M} \\
\texttt{/F} & Rest of line & \texttt{COMMAND/F} \\
\end{tabular}
\caption{ReadArgs template modifiers}
\end{table}

Modifiers can be combined: \texttt{COUNT/N/A} means required number, \texttt{OUTPUT/K/A} means required keyword.

\subsection{Argument Types}

\begin{lstlisting}
from amiga::workbench import Args

fn argument_types() {
    match Args::parse("FILE,COUNT/N,VERBOSE/S,FILES/M") {
        Result::Ok(args) => {
            // String argument
            match args.get_str("FILE") {
                Option::Some(f) => write("File: %s\n", f.as_cstr()),
                Option::None => println("No file specified"),
            }

            // Numeric argument
            let count = args.get_number("COUNT") or { 10 }  // Default 10
            write("Count: %ld\n", count)

            // Switch (boolean)
            if args.get_switch("VERBOSE") {
                println("Verbose mode enabled")
            }

            // Multiple arguments
            match args.get_multi("FILES") {
                Option::Some(var files) => {
                    while var i < files.len {
                        match files.get(i) {
                            Option::Some(f) => write("  %s\n", f.as_cstr()),
                            Option::None => {},
                        }
                        i++
                    }
                },
                Option::None => {},
            }
        },
        Result::Err(_) => {},
    }
}
\end{lstlisting}

\subsection{ArgParser Builder}

For programmatic template construction:

\begin{lstlisting}
from amiga::workbench import ArgParser, ArgType

fn builder_example() -> Result<(), ArgsError> {
    var parser = ArgParser::new()
    parser.add("FROM", ArgType::Required)
    parser.add("TO", ArgType::Required)
    parser.add("VERBOSE", ArgType::Switch)
    parser.add("BUFFER", ArgType::Number)

    let args = parser.parse()?

    // Retrieve values just like with Args::parse()
    let from = args.get_str("FROM") or { return Result::Err(ArgsError::MissingRequired) }
    let to = args.get_str("TO") or { return Result::Err(ArgsError::MissingRequired) }
    let verbose = args.get_switch("VERBOSE")
    let buffer_size = args.get_number("BUFFER") or { 4096 }  // Default

    if verbose {
        write("Copying %s to %s (buffer: %ld)\n", from.as_cstr(), to.as_cstr(), buffer_size)
    }

    return Result::Ok(())
}
\end{lstlisting}

\begin{table}[h]
\centering
\begin{tabular}{ll}
\textbf{ArgType} & \textbf{Template Equivalent} \\
\hline
\texttt{Optional} & (none) \\
\texttt{Required} & \texttt{/A} \\
\texttt{Switch} & \texttt{/S} \\
\texttt{Keyword} & \texttt{/K} \\
\texttt{Number} & \texttt{/N} \\
\texttt{Multi} & \texttt{/M} \\
\texttt{RestOfLine} & \texttt{/F} \\
\texttt{RequiredNumber} & \texttt{/N/A} \\
\texttt{RequiredKeyword} & \texttt{/K/A} \\
\end{tabular}
\caption{ArgType enum values}
\end{table}

\subsection{Argument Errors}

\begin{table}[h]
\centering
\begin{tabular}{ll}
\textbf{ArgsError} & \textbf{Meaning} \\
\hline
\texttt{Dos(DosError)} & Underlying DOS error (wrapped) \\
\texttt{InvalidTemplate} & Template syntax error \\
\texttt{MissingRequired} & Required argument not provided \\
\texttt{InvalidValue} & Argument value is invalid \\
\texttt{TooManyArgs} & Template has $>$ 32 arguments \\
\texttt{AllocationFailed} & Memory allocation failed \\
\texttt{ArgNotFound} & Argument name not in template \\
\end{tabular}
\caption{ArgsError enum variants}
\end{table}

%%%%%%%%%%%%%%%%%%%%%%%%%%%%%%%%%%%%%%%%%%%%%%%%%%%%%%%%%%%%%%%%%%%%%%%%%%%%%%%
\section{Workbench Support}
\label{sec:workbench}
%%%%%%%%%%%%%%%%%%%%%%%%%%%%%%%%%%%%%%%%%%%%%%%%%%%%%%%%%%%%%%%%%%%%%%%%%%%%%%%

Programs can be launched from either the CLI or Workbench. The two modes require different handling.

\subsection{Detecting Launch Mode}

\begin{lstlisting}
from amiga::workbench import is_workbench, Args, get_workbench_startup

pub fn main() -> i32 {
    if is_workbench() {
        // Launched from icon - no stdin/stdout
        return handle_workbench()
    } else {
        // Launched from CLI
        return handle_cli()
    }
}

fn handle_cli() -> i32 {
    match Args::parse("FILE/A") {
        Result::Ok(args) => {
            let file = args.get_str("FILE") or { return 10 }
            // Process file...
            return 0
        },
        Result::Err(_) => return 10,
    }
}

fn handle_workbench() -> i32 {
    match get_workbench_startup() {
        Option::Some(wb) => {
            // Process Workbench arguments
            if wb.num_args() > 1 {
                // User shift-clicked files
                let arg = wb.arg(1) or { return 20 }
                // Process first selected file...
            }
            return 0
        },
        Option::None => return 20,
    }
}
\end{lstlisting}

\subsection{Workbench Startup Message}

When launched from Workbench, the program receives a \texttt{WBStartup} message containing information about the icon clicked and any files selected:

\begin{notebox}
\textbf{Message reply is automatic.} The Novus runtime startup code retrieves the WBStartup message and handles replying to it when the program exits. You do \textbf{not} need to manually reply to the message---the runtime handles this for you.
\end{notebox}

\begin{lstlisting}
from amiga::workbench import get_workbench_startup, WorkbenchStartup

fn workbench_example() {
    let wb = get_workbench_startup() or { return }

    // Number of arguments (index 0 = program icon)
    let count = wb.num_args()
    write("Got %ld arguments\n", (i32)count)

    // Get program's own icon (always index 0)
    match wb.arg(0) {
        Option::Some(prog) => {
            match prog.name() {
                Option::Some(name) => {
                    write("Program: %s\n", name.as_cstr())
                },
                Option::None => {},
            }
        },
        Option::None => {},
    }

    // Process any selected files (indices 1..n)
    while var i: u32 = 1 {
        if i >= count { break }

        match wb.arg(i) {
            Option::Some(arg) => {
                match arg.full_path() {
                    Result::Ok(path) => {
                        write("File: %s\n", path.as_cstr())
                    },
                    Result::Err(_) => {},
                }
            },
            Option::None => {},
        }
        i++
    }
}
\end{lstlisting}

\subsection{WBArgument Structure}

Each Workbench argument provides:

\begin{lstlisting}
pub struct WBArgument {
    lock: i32,      // Lock on containing directory
    name_ptr: *u8,  // Filename (not full path)
}

impl WBArgument {
    /// Get directory lock (for CurrentDir)
    pub fn lock(&self) -> i32

    /// Get filename only
    pub fn name(&self) -> Option<Str>

    /// Get full path (returns String, auto-freed on drop)
    pub fn full_path(&self) -> Result<String, DosError>
}
\end{lstlisting}

\begin{tipbox}
The \texttt{lock} field is a DOS lock on the \emph{directory} containing the file, not the file itself. Use it with \texttt{CurrentDir()} to access the file by name, or call \texttt{full\_path()} to get the complete path. The \texttt{String} returned by \texttt{full\_path()} is automatically freed when it goes out of scope.
\end{tipbox}

\subsection{Tool Window}

If the program's icon has a \texttt{TOOLWINDOW} tooltype, you can get the window specification:

\begin{lstlisting}
fn check_tool_window(wb: &WorkbenchStartup) {
    match wb.tool_window() {
        Option::Some(window) => {
            // window is something like "CON:0/0/640/200/Output"
            write("Tool window: %s\n", window.as_cstr())
        },
        Option::None => {
            // No TOOLWINDOW specified
        },
    }
}
\end{lstlisting}

%%%%%%%%%%%%%%%%%%%%%%%%%%%%%%%%%%%%%%%%%%%%%%%%%%%%%%%%%%%%%%%%%%%%%%%%%%%%%%%
\section{Error Handling}
\label{sec:io-errors}
%%%%%%%%%%%%%%%%%%%%%%%%%%%%%%%%%%%%%%%%%%%%%%%%%%%%%%%%%%%%%%%%%%%%%%%%%%%%%%%

I/O operations return \texttt{Result} types for proper error handling:

\begin{lstlisting}
from amiga::sys::dos import OwnedFileHandle, DosError
from amiga::raw::dos import IoErr

fn robust_io() -> Result<(), DosError> {
    // The ? operator propagates errors
    let file = OwnedFileHandle::open("RAM:data.txt", MODE_OLDFILE)?

    var buffer: [u8; 1024] = [0; 1024]
    let result = file.read_all((*u8)&buffer[0], 1024)?

    return Result::Ok(())
}

fn detailed_error_handling() {
    match OwnedFileHandle::open("NONEXISTENT:", MODE_OLDFILE) {
        Result::Ok(file) => {
            // Use file...
        },
        Result::Err(e) => {
            write("Error: %s\n", e.message())

            // Get detailed DOS error code
            let code = IoErr()
            write("DOS error code: %ld\n", code)
        },
    }
}
\end{lstlisting}

\subsection{Common DOS Errors}

\begin{table}[h]
\centering
\begin{tabular}{lll}
\textbf{Error Code} & \textbf{Constant} & \textbf{Meaning} \\
\hline
103 & \texttt{ERROR\_NO\_FREE\_STORE} & Out of memory \\
205 & \texttt{ERROR\_OBJECT\_NOT\_FOUND} & File not found \\
212 & \texttt{ERROR\_OBJECT\_IN\_USE} & File locked by another \\
214 & \texttt{ERROR\_DISK\_FULL} & No space on disk \\
221 & \texttt{ERROR\_DISK\_WRITE\_PROTECTED} & Write-protected \\
225 & \texttt{ERROR\_SEEK\_ERROR} & Invalid seek position \\
232 & \texttt{ERROR\_NO\_MORE\_ENTRIES} & Directory scan complete \\
\end{tabular}
\caption{Common DOS error codes}
\end{table}

%%%%%%%%%%%%%%%%%%%%%%%%%%%%%%%%%%%%%%%%%%%%%%%%%%%%%%%%%%%%%%%%%%%%%%%%%%%%%%%
\section{I/O Patterns}
\label{sec:io-patterns}
%%%%%%%%%%%%%%%%%%%%%%%%%%%%%%%%%%%%%%%%%%%%%%%%%%%%%%%%%%%%%%%%%%%%%%%%%%%%%%%

\subsection{Pattern 1: Read Entire File}

\begin{lstlisting}
from amiga::sys::dos import OwnedFileHandle, file_info
from std::memory::block import MemoryBlock
from amiga::raw::consts import MODE_OLDFILE, MEMF_ANY, OFFSET_BEGINNING

fn read_entire_file(path: Str) -> Result<MemoryBlock, DosError> {
    // Get file size
    let info = file_info(path)?
    let size = (u32)info.size

    // Allocate buffer
    let buffer = MemoryBlock::try_alloc(size, MEMF_ANY)?

    // Read file
    let file = OwnedFileHandle::open(path, MODE_OLDFILE)?
    let read = file.read_all(buffer.ptr(), size)?

    return Result::Ok(buffer)
}
\end{lstlisting}

\subsection{Pattern 2: Write with Backup}

\begin{lstlisting}
from amiga::sys::dos import OwnedFileHandle, file_exists, delete_file
from amiga::raw::consts import MODE_NEWFILE

fn write_with_backup(path: Str, data: *u8, len: u32) -> Result<(), DosError> {
    // Create backup filename
    var backup_path = String::from(path)
    backup_path.push_str(".bak")

    // Delete old backup
    if file_exists(backup_path.as_str()) {
        let _ = delete_file(backup_path.as_str())
    }

    // Rename current to backup (if exists)
    if file_exists(path) {
        // Use Rename() from FFI...
    }

    // Write new file
    let file = OwnedFileHandle::open(path, MODE_NEWFILE)?
    file.write_all(data, len)?

    return Result::Ok(())
}
\end{lstlisting}

\subsection{Pattern 3: Process File Line by Line}

\begin{lstlisting}
fn process_lines(path: Str, handler: fn(*u8, u32) -> bool) -> Result<u32, DosError> {
    let file = OwnedFileHandle::open(path, MODE_OLDFILE)?

    var buffer: [u8; 1024] = [0; 1024]
    var line_count: u32 = 0

    loop {
        match file.read_line((*u8)&buffer[0], 1024) {
            Result::Ok(len) => {
                if len == 0 { break }  // EOF

                // Call handler, stop if it returns false
                if !handler((*u8)&buffer[0], len) {
                    break
                }
                line_count++
            },
            Result::Err(_) => break,
        }
    }

    return Result::Ok(line_count)
}
\end{lstlisting}

%%%%%%%%%%%%%%%%%%%%%%%%%%%%%%%%%%%%%%%%%%%%%%%%%%%%%%%%%%%%%%%%%%%%%%%%%%%%%%%
\section{I/O Best Practices}
\label{sec:io-best-practices}
%%%%%%%%%%%%%%%%%%%%%%%%%%%%%%%%%%%%%%%%%%%%%%%%%%%%%%%%%%%%%%%%%%%%%%%%%%%%%%%

\begin{enumerate}
    \item \textbf{Always use RAII wrappers}---\texttt{OwnedFileHandle} and \texttt{DirLockHandle} prevent leaks
    \item \textbf{Check for Workbench vs CLI}---handle both launch modes
    \item \textbf{Use \texttt{\%ld} for integers}---AmigaOS expects 32-bit values
    \item \textbf{Handle errors explicitly}---use \texttt{Result} and the \texttt{?} operator
    \item \textbf{Buffer reads and writes}---minimize system calls for performance
    \item \textbf{Close files promptly}---don't hold files open longer than necessary
    \item \textbf{Check disk space}---verify writes succeeded (\texttt{ERROR\_DISK\_FULL})
\end{enumerate}

\begin{warningbox}
\textbf{DOS has limited file handles!} Each process can have at most ~20 files open simultaneously. If you open many files, close them promptly or batch your operations.
\end{warningbox}

